\documentclass[11pt]{article}
\usepackage[margin=1in]{geometry}
\usepackage{amsmath,amssymb}
\usepackage{booktabs}
\usepackage{longtable}
\usepackage[hidelinks]{hyperref}

% Compact list spacing without enumitem.
\let\olditemize\itemize
\renewcommand{\itemize}{\olditemize\setlength{\itemsep}{0pt}%
  \setlength{\parskip}{0pt}\setlength{\topsep}{2pt}}

\title{Equation Reference for the 1D Charring-Ablator Solver\\
\large Every equation, every symbol, in plain terms}
\author{Grey Goodwin}
\date{}

\begin{document}
\maketitle

\section{Which Equations the Surrogate Uses}

The GNN imitates the solver's result. The solver computes the energy balance,
inverts for temperature, solves Darcy's law, and marches. The surrogate skips all
of that and predicts the next state directly from the current one, using a single
bounded stencil it learned from watching the solver's output.

So the surrogate does not solve this system. It reproduces part of it exactly,
replaces the transport with a learned bounded update, and never computes the rest
at all. The table below places every equation in one of those three groups.

\begin{longtable}{@{}p{0.07\textwidth} p{0.38\textwidth} p{0.48\textwidth}@{}}
\toprule
\textbf{Eq.} & \textbf{What it is} & \textbf{Role in the surrogate} \\
\midrule
\endfirsthead
\toprule
\textbf{Eq.} & \textbf{What it is} & \textbf{Role in the surrogate} \\
\midrule
\endhead
\multicolumn{3}{@{}l}{\textit{Computed in the forward pass, unchanged}} \\
15 & Decomposition state & Virgin fraction $W$, recomputed every sub-step \\
16 & Arrhenius rate law & The solver's own $A$, $E$, $\psi$, $T_{\min}$, held fixed \\
17 & Closed-form integration & The exact update, evaluated under \texttt{no\_grad} \\
18 & Bulk density, gas generation & Supplies $\dot m$ to the learned coupling \\
\midrule
\multicolumn{3}{@{}l}{\textit{Replaced by the learned bounded update}} \\
2, 3 & Flux coefficients, face weights & Absorbed into the single rate $D$ \\
4 & Total cell enthalpy & Not formed. The model predicts $T$ directly \\
5 & The energy update & Replaced by $u_i \leftarrow u_i + \sum_j D(u_j-u_i) + S_i$ \\
6 & Conduction fluxes & Replaced by $D\,(u_j - u_i)$ \\
7 & Advected gas enthalpy & $\dot m$ becomes a diffusive channel, not an advective one \\
8 & Chemical energy term & No energy bookkeeping. A learned coupling $c_T \dot m$ instead \\
9 & The continuous form & The operator the surrogate learns to imitate \\
10, 11 & Face state, face conductivity & One learned rate replaces the mixed conductivity \\
12, 13 & Temperature predictor and Newton correction & Unnecessary. $T$ is a predicted channel \\
14 & The time step & No CFL limit. The convex update is stable at any step \\
27, 28 & Darcy's law, discrete Darcy & Replaced by a conductance channel on $p_g$ \\
30, 31, 32 & Gas mass conservation & Not enforced. The weakest substitution in the model \\
\midrule
\multicolumn{3}{@{}l}{\textit{Upstream only: shapes the training data or arrives as input}} \\
1 & Mesh geometry & $\Delta x$ enters as an edge feature \\
19, 20 & Inert density, char progress $\tau$ & Solver only. $\rho$ is an input feature \\
21, 22 & Char fraction $\beta$, porosity & Computed in the data pipeline as input features \\
23 & Property mixing & Solver only. The surrogate has no material properties \\
24, 25 & Formation and sensible enthalpy & Solver only. No enthalpy in the surrogate \\
26 & Emissivity & Solver only, through the surface balance \\
29 & Permeability blending & Solver only \\
33 & Gas properties & Solver only. $p_g$ arrives as a state channel \\
34, 35, 36 & Surface balance, blowing, B-prime & Solver only. Their effect reaches the model through the ghost cell \\
37, 38 & Ghost cells & The forcing. Re-imposed from truth at every step \\
\bottomrule
\end{longtable}

\noindent
\textbf{The short version.} Four equations survive intact, all of them chemistry.
Eighteen are replaced by a single bounded stencil carrying ten learned
parameters. The remaining sixteen never appear in the surrogate at all. They
shape the data it trains on, or they arrive as boundary forcing it does not have
to compute.

\bigskip

\noindent
\textbf{What this costs.} The gas equations are the honest weak point. Equations
27 through 32 describe pressure-driven flow through a porous solid, and the
surrogate substitutes diffusion on $p_g$ and $\dot m$. That is a modeling
convenience rather than a discretization of Darcy's law, and the gas channels
carry no physical reading as a result.

\clearpage
\section{How a Heat Shield Protects}

The material protects the vehicle by being consumed. That happens through four
separate mechanisms, and every one of them appears somewhere in the equations
that follow. Keeping them distinct is worth the effort, because they are easy to
collapse into a single vague statement about ablation.

\subsection*{The char insulates}

As the resin decomposes it leaves behind a porous carbon skeleton. That char has
a much lower thermal conductivity than the virgin material, so it slows the heat
reaching the material still intact behind it. This enters through the property
mixing of Eq.~(23), where conductivity is blended between virgin and char values
by the char-progress variable $\tau$ of Eq.~(20).

\subsection*{The char re-radiates}

The outer surface of the char reaches well over 1500~K, and a surface that hot
radiates energy back out. The term is $\varepsilon \sigma T_w^4$ in the surface
energy balance of Eq.~(34), with the emissivity itself blended by char state in
Eq.~(26). This is a large term at high heating rates, because radiation scales
with the fourth power of temperature.

\subsection*{The reaction absorbs energy}

Breaking the chemical bonds in the resin takes energy, and that is energy which
does not go into raising the temperature. The solver captures this without a
separate heat-of-pyrolysis parameter. Virgin material and char have different
formation enthalpies, Eq.~(24), so the shift in baseline energy as the material
chars is exactly the reaction energy. It appears as $E^{\mathrm{chem}}$ in
Eq.~(8).

\subsection*{The pyrolysis gas cools and blocks}

The gas released by decomposition has to escape, and it percolates outward
through the hot char. This does two separate things.

On the way out, the gas picks up heat from the solid it passes through and
carries that energy toward the surface. This is the advection term $H_i = \dot m_i
h_g$ of Eq.~(7), driven by the Darcy flow of Eqs.~(27) and~(28).

At the surface, the escaping gas thickens the boundary layer and reduces how much
heat the external flow can deliver in the first place. This is the blowing
correction of Eq.~(35), which scales down the convective coefficient as a
function of the blowing parameter. It is a reduction in the incoming heat rather
than a removal of heat already absorbed.

\subsection*{A caution on ranking them}

Which mechanism dominates depends on the material and the heating rate, and the
balance shifts across a trajectory. This document does not rank them, and any
claim about their relative size should be supported by a source or by a term-wise
energy accounting of a specific case.

\clearpage
\section{Nomenclature}

\begin{longtable}{@{}p{0.13\textwidth} p{0.53\textwidth} p{0.14\textwidth} p{0.10\textwidth}@{}}
\toprule
\textbf{Symbol} & \textbf{Meaning} & \textbf{Units} & \textbf{First in} \\
\midrule
\endfirsthead
\toprule
\textbf{Symbol} & \textbf{Meaning} & \textbf{Units} & \textbf{First in} \\
\midrule
\endhead
\multicolumn{4}{@{}l}{\textit{Geometry and indexing}} \\
$L_0$ & Slab thickness & m & Eq.~1 \\
$N$ & Number of interior cells & -- & Eq.~1 \\
$N_s$ & Number of species & -- & Eq.~18 \\
$i$ & Cell or face index & -- & Eq.~1 \\
$k$ & Species index (also thermal conductivity) & -- & Eq.~15 \\
$n$ & Time-step index & -- & Eq.~5 \\
$m$ & Geometry selector (also Newton iteration index) & -- & Eq.~1 \\
$x$ & Position through the thickness & m & Eq.~9 \\
$t$ & Time & s & Eq.~9 \\
$\Delta x$ & Cell width & m & Eq.~1 \\
$\Delta t$ & Time step & s & Eq.~5 \\
$dA$ & Face area & m$^2$ & Eq.~1 \\
$dV$ & Cell volume & m$^3$ & Eq.~1 \\
$A$ & Cross-sectional area (continuous) & m$^2$ & Eq.~9 \\
$\mathrm{FC}^\pm$ & Flux coefficient & 1/m & Eq.~2 \\
$\mathrm{g}^\pm$ & Face interpolation weight & -- & Eq.~3 \\
$\mathrm{CFL}$ & Time-step safety divisor & -- & Eq.~14 \\
\midrule
\multicolumn{4}{@{}l}{\textit{Thermal state}} \\
$T$ & Temperature & K & Eq.~4 \\
$T_f$ & Temperature at a face & K & Eq.~7 \\
$T_w$ & Wall temperature & K & Eq.~34 \\
$T_{\mathrm{ext}}$ & External radiation sink temperature & K & Eq.~34 \\
$T_{\mathrm{ref}}$ & Reference temperature (298.15 K) & K & Eq.~24 \\
$E$ & Total energy in a cell & J & Eq.~4 \\
$E^{\mathrm{chem}}$ & Energy change from decomposition & J & Eq.~5 \\
$F^\pm$ & Conduction flux at a face & W/m$^2$ & Eq.~5 \\
$H$ & Gas enthalpy flow across a face & W & Eq.~5 \\
$q_w$ & Net heat flux into the wall & W/m$^2$ & Eq.~34 \\
$q_{\mathrm{rad}}$ & Incident radiative flux & W/m$^2$ & Eq.~34 \\
$\alpha$ & Inverse cell heat capacity & K/J & Eq.~12 \\
\midrule
\multicolumn{4}{@{}l}{\textit{Solid material}} \\
$\rho$ & Bulk solid density & kg/m$^3$ & Eq.~4 \\
$\rho_f$ & Density interpolated to a face & kg/m$^3$ & Eq.~10 \\
$\rho_{i,k}$ & Density of species $k$ & kg/m$^3$ & Eq.~15 \\
$\rho_{v,k}$, $\rho_{c,k}$ & Intrinsic virgin, char density of species $k$ & kg/m$^3$ & Eq.~15 \\
$\rho_{v,\mathrm{bulk}}$ & Bulk virgin density & kg/m$^3$ & Eq.~14 \\
$\rho_{c,\mathrm{bulk}}$ & Bulk char density & kg/m$^3$ & Eq.~21 \\
$\rho_{\mathrm{inert}}$ & Density of non-reacting species & kg/m$^3$ & Eq.~19 \\
$\rho_r,\rho_{vr},\rho_{cr}$ & Densities with inert part removed & kg/m$^3$ & Eq.~20 \\
$\tau$ & Char-progress variable (1 virgin, 0 char) & -- & Eq.~4 \\
$\beta$ & Linear char fraction (0 virgin, 1 char) & -- & Eq.~21 \\
$\varphi$ & Porosity & -- & Eq.~4 \\
$\varphi_v$, $\varphi_c$ & Virgin, char porosity & -- & Eq.~18 \\
$k_{\mathrm{mix}}$ & Blended thermal conductivity & W/m/K & Eq.~9 \\
$k_v$, $k_c$ & Virgin, char thermal conductivity & W/m/K & Eq.~11 \\
$k_f$ & Conductivity at a face & W/m/K & Eq.~6 \\
$c_{p,\mathrm{mix}}$ & Blended specific heat & J/kg/K & Eq.~12 \\
$c_{p,v}$, $c_{p,c}$ & Virgin, char specific heat & J/kg/K & Eq.~23 \\
$h_s$ & Sensible enthalpy of the solid & J/kg & Eq.~4 \\
$h_{s,\mathrm{form}}$ & Formation enthalpy of the solid & J/kg & Eq.~8 \\
$e_v$, $e_c$ & Integrated heat capacity, virgin and char & J/kg & Eq.~24 \\
$\varepsilon$ & Emissivity & -- & Eq.~26 \\
$\varepsilon_v$, $\varepsilon_c$ & Virgin, char emissivity & -- & Eq.~26 \\
\midrule
\multicolumn{4}{@{}l}{\textit{Pyrolysis kinetics}} \\
$W_k$ & Normalized reactive mass fraction & -- & Eq.~15 \\
$A_k$ & Pre-exponential factor & 1/s & Eq.~16 \\
$\psi_k$ & Reaction order & -- & Eq.~16 \\
$E_k$ & Activation temperature ($E/R$) & K & Eq.~16 \\
$\gamma_k$ & Species mass-fraction weight & -- & Eq.~18 \\
$\dot\omega_g$ & Volumetric gas generation rate & kg/m$^3$/s & Eq.~18 \\
$\dot\omega_{\mathrm{net}}$ & Gas source net of storage & kg/m$^3$/s & Eq.~31 \\
$\mathcal{I}$ & Set of inert species & -- & Eq.~19 \\
\midrule
\multicolumn{4}{@{}l}{\textit{Pyrolysis gas}} \\
$\rho_g$ & Gas density & kg/m$^3$ & Eq.~4 \\
$p_g$ & Gas pressure & Pa & Eq.~4 \\
$p_{\mathrm{amb}}$ & Ambient pressure at the surface & Pa & Eq.~32 \\
$\dot m$ & Gas mass flow across a face & kg/s & Eq.~7 \\
$\dot m''$ & Gas mass flux per unit area & kg/m$^2$/s & Eq.~9 \\
$u_g$ & Gas specific internal energy & J/kg & Eq.~4 \\
$h_g$ & Gas enthalpy relative to $T_{\mathrm{ref}}$ & J/kg & Eq.~7 \\
$h_g^{\mathrm{abs}}$ & Absolute gas enthalpy & J/kg & Eq.~34 \\
$c_{p,g}$, $c_{v,g}$ & Gas specific heats & J/kg/K & Eq.~12 \\
$R$ & Specific gas constant (tabulated) & J/kg/K & Eq.~33 \\
$\mu$ & Gas dynamic viscosity & Pa$\cdot$s & Eq.~27 \\
$K_0$ & Permeability & m$^2$ & Eq.~27 \\
$K_{0,v}$, $K_{0,c}$ & Virgin, char permeability & m$^2$ & Eq.~29 \\
$C$ & Gas conductance of a face & kg/s/Pa & Eq.~28 \\
\midrule
\multicolumn{4}{@{}l}{\textit{Aerothermal boundary}} \\
$\rho_e u_e C_H$ & Heat-transfer coefficient & kg/m$^2$/s & Eq.~34 \\
$h_r$ & Recovery enthalpy & J/kg & Eq.~34 \\
$h_w$ & Wall enthalpy & J/kg & Eq.~34 \\
$\Phi$ & Blowing correction factor & -- & Eq.~34 \\
$\varphi_b$ & Dimensionless blowing parameter & -- & Eq.~35 \\
$\lambda$ & Blowing coefficient & -- & Eq.~35 \\
$B'_g$ & Dimensionless gas blowing rate & -- & Eq.~36 \\
$\sigma$ & Stefan-Boltzmann constant & W/m$^2$/K$^4$ & Eq.~34 \\
\bottomrule
\end{longtable}

\section*{Symbols that are reused}

Three letters carry more than one meaning. They are distinguished by context, but
they are worth flagging.

\begin{itemize}
  \item $m$ is the geometry selector in Eq.~(1) and the Newton iteration counter
        in Eq.~(13).
  \item $k$ is thermal conductivity throughout the energy equation and the
        species index throughout the pyrolysis section.
  \item $\varphi$ is porosity in Eq.~(22) and, with a subscript $b$, the
        dimensionless blowing parameter in Eq.~(35).
\end{itemize}

\noindent
Two further pairs look similar but are not the same quantity.

\begin{itemize}
  \item $\dot m$ is a mass flow rate across a face [kg/s], while $\dot m''$ is a
        mass flux per unit area [kg/m$^2$/s].
  \item $h_g$ is gas enthalpy referenced to $T_{\mathrm{ref}}$, while
        $h_g^{\mathrm{abs}}$ is absolute gas enthalpy. Only the absolute form
        appears in the surface balance.
\end{itemize}

\clearpage
\noindent
This is a working reference, not a manuscript. Every equation implemented by the
reference solver is listed in the order it appears in the chapter draft. Under
each one is a one-sentence statement of what it does and a list identifying every
symbol in it. A symbol introduced earlier is marked as such rather than being
redefined. The consolidated nomenclature is given ahead of this section.

\vspace{1em}
\hrule
\vspace{1em}

\section{Domain and Discretization}

\subsection*{Equation 1: mesh geometry}
\begin{equation}
\Delta x = \frac{L_0}{N}, \qquad
dA_i = (2x_i)^{m}\,\pi^{m(3-m)/2}, \qquad
dV_i = \frac{2^{m}\pi^{m(3-m)/2}}{m+1}\left(x_{i+1}^{m+1} - x_i^{m+1}\right)
\end{equation}
\textbf{What it is.} Bookkeeping that chops the slab into control volumes. It
defines how wide each cell is, how much area its faces have, and how much volume
it contains.
\begin{itemize}
  \item $\Delta x$ = width of one cell [m].
  \item $L_0$ = total thickness of the material slab [m].
  \item $N$ = number of interior cells.
  \item $dA_i$ = area of the face on the left side of cell $i$ [m$^2$].
  \item $dV_i$ = volume of cell $i$ [m$^3$].
  \item $x_i$ = position of face $i$ [m].
  \item $m$ = geometry selector. $0$ = flat slab, $1$ = cylinder, $2$ = sphere.
  \item $i$ = cell index, running $1$ to $N$.
\end{itemize}
\textit{Note.} All cases here use $m=0$, which collapses this to $dA_i = 1$ and
$dV_i = \Delta x$.

\subsection*{Equation 2: flux coefficients}
\begin{equation}
\mathrm{FC}^-_i = \frac{2}{\Delta x_i + \Delta x_{i-1}}, \qquad
\mathrm{FC}^+_i = \frac{2}{\Delta x_{i+1} + \Delta x_i}
\end{equation}
\textbf{What it is.} The geometric part of a finite-difference derivative. It is
one over the center-to-center distance between two neighboring cells, so
multiplying it by a temperature difference gives a temperature gradient.
\begin{itemize}
  \item $\mathrm{FC}^-_i$ = coefficient for the left face of cell $i$ [1/m].
  \item $\mathrm{FC}^+_i$ = coefficient for the right face of cell $i$ [1/m].
  \item $\Delta x$, $i$ = as in Eq.~(1).
\end{itemize}

\subsection*{Equation 3: face interpolation weights}
\begin{equation}
\mathrm{g}^-_i = \frac{\Delta x_{i-1}}{\Delta x_{i-1} + \Delta x_i}, \qquad
\mathrm{g}^+_i = \frac{\Delta x_{i+1}}{\Delta x_i + \Delta x_{i+1}}
\end{equation}
\textbf{What it is.} A weighting fraction between $0$ and $1$ used to estimate a
value at a face from the values at the two cells on either side. On a uniform
mesh both equal $1/2$, which is a plain average.
\begin{itemize}
  \item $\mathrm{g}^-_i$, $\mathrm{g}^+_i$ = interpolation weights, dimensionless.
  \item $\Delta x$, $i$ = as in Eq.~(1).
\end{itemize}

\section{Energy Equation}

\subsection*{Equation 4: the conserved quantity}
\begin{equation}
E_i = \underbrace{\rho_i\, h_s(T_i, \tau_i)\, dV_i}_{\text{solid}}
    + \underbrace{\varphi_i\, \rho_g(T_i, p_{g,i})\, u_g(T_i)\, dV_i}_{\text{gas}}
\end{equation}
\textbf{What it is.} The total energy stored in one cell. This is the quantity
the solver actually tracks, instead of temperature. It is the energy in the solid
plus the energy in the gas sitting in the pores.
\begin{itemize}
  \item $E_i$ = total energy in cell $i$ [J].
  \item $\rho_i$ = bulk density of the solid in cell $i$ [kg/m$^3$].
  \item $h_s$ = sensible enthalpy of the solid per unit mass [J/kg]. Defined in
        Eq.~(25).
  \item $T_i$ = temperature of cell $i$ [K].
  \item $\tau_i$ = char-progress variable, $1$ = virgin, $0$ = fully charred.
        Defined in Eq.~(20).
  \item $\varphi_i$ = porosity, the fraction of the cell that is void space.
        Defined in Eq.~(22).
  \item $\rho_g$ = density of the pyrolysis gas [kg/m$^3$]. Defined in Eq.~(33).
  \item $p_{g,i}$ = gas pressure in cell $i$ [Pa].
  \item $u_g$ = specific internal energy of the gas [J/kg].
  \item $dV_i$ = as in Eq.~(1).
\end{itemize}

\subsection*{Equation 5: the energy update}
\begin{equation}
E_i^{n+1} = E_i^{n}
  + \Delta t \Big[ \underbrace{-F^-_i\, dA_i + F^+_i\, dA_{i+1}}_{\text{conduction}}
  - \underbrace{\left(H_{i+1} - H_i\right)}_{\text{gas advection}} \Big]
  + E^{\mathrm{chem}}_i
\end{equation}
\textbf{What it is.} The central equation of the solver. It says the energy in a
cell after one time step equals what it had before, plus heat conducted in across
its two faces, minus energy carried away by gas flowing through, plus energy
released or absorbed by the chemical decomposition.
\begin{itemize}
  \item $n$, $n+1$ = time-step indices, meaning before and after the step.
  \item $\Delta t$ = size of the time step [s]. Set by Eq.~(14).
  \item $F^-_i$, $F^+_i$ = conduction heat flux at the left and right faces
        [W/m$^2$]. Defined in Eq.~(6).
  \item $H_i$ = enthalpy carried by gas across face $i$ [W]. Defined in Eq.~(7).
  \item $E^{\mathrm{chem}}_i$ = energy change from chemical decomposition [J].
        Defined in Eq.~(8).
  \item $E_i$, $dA_i$ = as in Eqs.~(4) and~(1).
\end{itemize}

\subsection*{Equation 6: conduction fluxes}
\begin{equation}
F^-_i = k_{f,i}\,\mathrm{FC}^-_i\,(T_i - T_{i-1}), \qquad
F^+_i = k_{f,i+1}\,\mathrm{FC}^+_i\,(T_{i+1} - T_i)
\end{equation}
\textbf{What it is.} Fourier's law of conduction written on the mesh. Heat flows
down the temperature gradient at a rate set by the conductivity.
\begin{itemize}
  \item $k_{f,i}$ = thermal conductivity evaluated at face $i$ [W/m/K]. Defined
        in Eq.~(11).
  \item $F^-_i$, $F^+_i$, $\mathrm{FC}^\pm_i$, $T_i$ = as in Eqs.~(5), (2), (4).
\end{itemize}

\subsection*{Equation 7: advected gas enthalpy}
\begin{equation}
H_i = \dot m_i\, h_g(T_{f,i}), \qquad
T_{f,i} = T_{i-1} + (T_i - T_{i-1})\,\mathrm{g}^-_i
\end{equation}
\textbf{What it is.} Hot gas moving through the material carries energy with it.
This is how much, equal to the mass flow rate times the enthalpy per unit mass.
\begin{itemize}
  \item $\dot m_i$ = mass flow rate of gas across face $i$ [kg/s]. Solved in
        Eq.~(27).
  \item $h_g$ = enthalpy of the gas per unit mass, measured relative to the
        reference temperature [J/kg].
  \item $T_{f,i}$ = temperature interpolated to face $i$ [K].
  \item $H_i$, $\mathrm{g}^-_i$, $T_i$ = as in Eqs.~(5), (3), (4).
\end{itemize}

\subsection*{Equation 8: the chemical energy term}
\begin{equation}
E^{\mathrm{chem}}_i = \Big[\rho_i^{n+1} h_{s,\mathrm{form}}(\tau_i^{n+1})
  - \rho_i^{n} h_{s,\mathrm{form}}(\tau_i^{n})\Big] dV_i
\end{equation}
\textbf{What it is.} The heat of pyrolysis, expressed as a bookkeeping change
rather than a source term. As the material chars, its baseline enthalpy shifts,
and that shift is the energy the reaction consumed or released.
\begin{itemize}
  \item $h_{s,\mathrm{form}}$ = formation enthalpy of the solid at the reference
        temperature, blended by char state [J/kg]. Defined in Eq.~(24).
  \item $E^{\mathrm{chem}}_i$, $\rho_i$, $\tau_i$, $dV_i$, $n$ = as in
        Eqs.~(5), (4), (1).
\end{itemize}
\textit{Note.} This is the only place pyrolysis energy enters. There is no
separate heat-of-pyrolysis parameter anywhere in the solver.

\subsection*{Equation 9: the continuous form}
\begin{equation}
\frac{\partial}{\partial t}\Big[\rho\,h_s(T,\tau) + \varphi\,\rho_g u_g\Big]
= \frac{1}{dV}\frac{\partial}{\partial x}\!\left(A\,k_{\mathrm{mix}}\frac{\partial T}{\partial x}\right)
- \frac{\partial}{\partial x}\!\left(\dot m''\, h_g\right)
- \frac{\partial}{\partial t}\Big[\rho\,h_{s,\mathrm{form}}(\tau)\Big]
\end{equation}
\textbf{What it is.} Equations (4) through (8) written as a partial differential
equation instead of a discrete update. This is the governing equation the solver
approximates. It is the charring-ablator energy equation.
\begin{itemize}
  \item $t$ = time [s].
  \item $x$ = position through the thickness [m].
  \item $A$ = cross-sectional area [m$^2$], the continuous version of $dA$.
  \item $k_{\mathrm{mix}}$ = char-blended thermal conductivity [W/m/K]. Defined
        in Eq.~(23).
  \item $\dot m''$ = gas mass flux per unit area [kg/m$^2$/s]. Note the double
        prime distinguishes it from $\dot m$ in Eq.~(7), which is per face and
        not per area.
  \item All other symbols as previously defined.
\end{itemize}

\subsection*{Equation 10: face state interpolation}
\begin{equation}
\rho_{f,i} = \tfrac12\left(\rho_{i-1} + \rho_i\right), \qquad
T_{f,i} = T_i + (T_{i-1} - T_i)\,\mathrm{g}^+_{i-1}
\end{equation}
\textbf{What it is.} Estimates the density and temperature at a face before the
conductivity is computed there.
\begin{itemize}
  \item $\rho_{f,i}$ = density interpolated to face $i$ [kg/m$^3$].
  \item $T_{f,i}$, $\rho_i$, $\mathrm{g}^+_i$ = as in Eqs.~(7), (4), (3).
\end{itemize}

\subsection*{Equation 11: face conductivity}
\begin{equation}
k_{f,i} = \tau_{f,i}\, k_v(T_{f,i}) + \left(1 - \tau_{f,i}\right) k_c(T_{f,i}),
\qquad \tau_{f,i} = \tau(\rho_{f,i})
\end{equation}
\textbf{What it is.} The conductivity used in Eq.~(6). Virgin and char material
conduct heat differently, so the two tabulated curves are mixed according to how
charred the face is.
\begin{itemize}
  \item $k_v(T)$ = tabulated conductivity of virgin material [W/m/K].
  \item $k_c(T)$ = tabulated conductivity of char [W/m/K].
  \item $\tau_{f,i}$ = char-progress variable evaluated at the face.
  \item $k_{f,i}$, $T_{f,i}$, $\rho_{f,i}$ = as in Eqs.~(6), (7), (10).
\end{itemize}
\textit{Note.} The conductivity function is evaluated once at the interpolated
face state. It is not an average of the two neighboring cell conductivities. The
density used is the one from the start of the step, before pyrolysis.

\subsection*{Equation 12: temperature predictor}
\begin{equation}
T^{(0)}_i = T^n_i + \alpha_i\left(\text{net flux} - \text{advection}\right),
\qquad
\alpha_i = \frac{\Delta t / dV_i}{\rho^{n+1}_i c_{p,\mathrm{mix}}(T^n_i, \bar\tau_i)
  + \bar\varphi_i \rho_g c_{v,g}}
\end{equation}
\textbf{What it is.} A first guess at the new temperature, obtained by dividing
the energy change by the heat capacity. It starts the iteration of Eq.~(13).
\begin{itemize}
  \item $T^{(0)}_i$ = initial guess for the new temperature [K]. The parenthesized
        superscript counts iterations, not time steps.
  \item $\alpha_i$ = inverse heat capacity of the cell [K/J].
  \item $c_{p,\mathrm{mix}}$ = char-blended specific heat [J/kg/K]. Defined in
        Eq.~(23).
  \item $c_{v,g}$ = specific heat of the gas at constant volume [J/kg/K].
        Defined in Eq.~(33).
  \item $\bar\tau_i$, $\bar\varphi_i$ = char variable and porosity averaged
        between the old and new step.
  \item All other symbols as previously defined.
\end{itemize}

\subsection*{Equation 13: Newton correction for temperature}
\begin{equation}
T^{(m+1)}_i = T^{(m)}_i +
\frac{E^{n+1}_i - \left[\rho^{n+1}_i h_s(T^{(m)}_i, \tau^{n+1}_i)
  + \varphi^{n+1}_i \rho_g u_g\right] dV_i}
{\left[\rho^{n+1}_i c_{p,\mathrm{mix}}(T^{(m)}_i, \tau^{n+1}_i)
  + \varphi^{n+1}_i \rho_g c_{v,g}\right] dV_i}
\end{equation}
\textbf{What it is.} Recovers temperature from energy. Since enthalpy depends
nonlinearly on temperature, the solver guesses a temperature, computes what
energy that implies, and corrects by the mismatch divided by the heat capacity.
\begin{itemize}
  \item $m$ = iteration counter. Note this reuses the letter $m$ from Eq.~(1),
        where it meant the geometry selector. They are unrelated.
  \item All other symbols as previously defined.
\end{itemize}
\textit{Note.} Applied exactly three times with no convergence test.

\subsection*{Equation 14: the time step}
\begin{equation}
\Delta t = \frac{1}{\mathrm{CFL}} \cdot
\frac{\Delta x_{\max}^2\, \rho_{v,\mathrm{bulk}}\, c_{p,\mathrm{mix}}(600\,\mathrm{K}, \tau_{\mathrm{init}})}
{2\, k_{\mathrm{mix}}(600\,\mathrm{K}, \tau_{\mathrm{init}})}
\end{equation}
\textbf{What it is.} The stability limit for an explicit diffusion scheme. Because
it scales with $\Delta x^2$, doubling the mesh resolution quarters the time step
and quadruples the number of steps. This is the cost the surrogate is built to
avoid.
\begin{itemize}
  \item $\mathrm{CFL}$ = a user safety divisor, dimensionless. It sits in the
        denominator, so a larger value gives a \emph{smaller} step. It is not a
        Courant number.
  \item $\Delta x_{\max}$ = the largest cell width in the mesh [m].
  \item $\rho_{v,\mathrm{bulk}}$ = bulk density of fully virgin material
        [kg/m$^3$].
  \item $\tau_{\mathrm{init}}$ = the initial char state, fully virgin.
  \item $600$~K = a hard-coded evaluation temperature for the properties.
  \item $c_{p,\mathrm{mix}}$, $k_{\mathrm{mix}}$, $\Delta t$ = as in Eqs.~(12),
        (9), (5).
\end{itemize}

\section{Pyrolysis}

\subsection*{Equation 15: decomposition state}
\begin{equation}
W_k = \mathrm{clip}\!\left(\frac{\rho_{i,k} - \rho_{c,k}}{\rho_{v,k}},\, 0,\, 1\right)
\end{equation}
\textbf{What it is.} How much of species $k$ is left to react, normalized to run
from $1$ (untouched) down to $0$ (fully decomposed).
\begin{itemize}
  \item $W_k$ = normalized reactive mass fraction of species $k$, dimensionless.
  \item $k$ = species index. Note this reuses the letter $k$ from Eq.~(6), where
        it meant thermal conductivity. Context distinguishes them.
  \item $\rho_{i,k}$ = current density of species $k$ [kg/m$^3$].
  \item $\rho_{v,k}$ = intrinsic virgin density of species $k$ [kg/m$^3$].
  \item $\rho_{c,k}$ = intrinsic char density of species $k$ [kg/m$^3$].
  \item $\mathrm{clip}(a,b,c)$ = clamps $a$ to lie between $b$ and $c$.
\end{itemize}

\subsection*{Equation 16: the Arrhenius rate law}
\begin{equation}
\frac{dW_k}{dt} = -A_k\, W_k^{\psi_k}\, \exp\!\left(-\frac{E_k}{T}\right)
\end{equation}
\textbf{What it is.} The chemistry. Decomposition speeds up exponentially with
temperature and slows as the reactive material is used up.
\begin{itemize}
  \item $A_k$ = pre-exponential factor, sets the overall rate [1/s].
  \item $\psi_k$ = reaction order, dimensionless. Controls how the rate falls off
        as material is consumed.
  \item $E_k$ = activation temperature [K]. Already divided by the gas constant,
        which is why no gas constant appears here.
  \item $W_k$, $T$, $t$ = as in Eqs.~(15), (4), (9).
\end{itemize}

\subsection*{Equation 17: closed-form integration}
\begin{equation}
W^{n+1}_k =
\begin{cases}
W_k \exp\!\left[-A_k \Delta t\, e^{-E_k/T}\right], & \psi_k = 1, \\[8pt]
\left[\max\!\left(0,\; W_k^{1-\psi_k} - \Delta t\,(1-\psi_k) A_k e^{-E_k/T}\right)\right]^{\frac{1}{1-\psi_k}},
& \psi_k \neq 1
\end{cases}
\end{equation}
\textbf{What it is.} The exact analytic solution of Eq.~(16) over one time step,
holding temperature fixed. The solver does not integrate the rate law
numerically. The two cases exist because the integral has a different form when
the reaction order is exactly one.
\begin{itemize}
  \item All symbols as in Eqs.~(15) and~(16).
\end{itemize}
\textit{Note.} There is no minimum reaction temperature. One is present in the
material file but never applied.

\subsection*{Equation 18: bulk density and gas generation}
\begin{equation}
\rho = \sum_{k=1}^{N_s} (1-\varphi_v)\,\gamma_k\, \rho_{i,k},
\qquad
\dot\omega_g = \sum_{k=1}^{N_s} (1-\varphi_v)\,\gamma_k\,
\frac{\max\!\left(0,\, \rho^{n}_{i,k} - \rho^{n+1}_{i,k}\right)}{\Delta t}
\end{equation}
\textbf{What it is.} Adds up the species to get the total solid density, and
converts the mass the solid lost into the rate at which pyrolysis gas is being
produced. Mass lost by the solid becomes gas.
\begin{itemize}
  \item $N_s$ = number of species. Three for the material used here.
  \item $\gamma_k$ = mass-fraction weight of species $k$, dimensionless.
  \item $\varphi_v$ = porosity of virgin material, dimensionless.
  \item $\dot\omega_g$ = volumetric gas generation rate [kg/m$^3$/s].
  \item $\rho$, $\rho_{i,k}$, $\Delta t$ = as in Eqs.~(4), (15), (5).
\end{itemize}

\section{Char State and Material Properties}

\subsection*{Equation 19: inert density}
\begin{equation}
\rho_{\mathrm{inert}} = \rho_{v,\mathrm{bulk}}
\frac{\sum_{k \in \mathcal{I}} \rho_{v,k}}{\sum_{k} \rho_{v,k}},
\qquad
\mathcal{I} = \left\{k : \rho_{v,k} = \rho_{c,k}\right\}
\end{equation}
\textbf{What it is.} Identifies the part of the material that never reacts and
measures its density, so it can be subtracted out before char progress is
computed. Without this, the inert fraction would compress the range of $\tau$.
\begin{itemize}
  \item $\rho_{\mathrm{inert}}$ = bulk density contributed by inert species
        [kg/m$^3$].
  \item $\mathcal{I}$ = the set of species whose virgin and char densities are
        equal, meaning they do not decompose.
  \item $\rho_{v,\mathrm{bulk}}$, $\rho_{v,k}$, $\rho_{c,k}$ = as in Eqs.~(14),
        (15).
\end{itemize}

\subsection*{Equation 20: char-progress variable}
\begin{equation}
\tau = \mathrm{clip}\!\left[
\frac{1 - \rho_{cr}/\rho_r}{1 - \rho_{cr}/\rho_{vr}},\, 0,\, 1\right]
\end{equation}
\textbf{What it is.} The measure of how charred the material is, used to blend
the thermal properties. It equals $1$ in virgin material and $0$ in fully charred
material. Note it is nonlinear in density.
\begin{itemize}
  \item $\tau$ = char-progress variable, dimensionless. First used in Eq.~(4).
  \item $\rho_r$ = current bulk density with the inert part removed [kg/m$^3$].
  \item $\rho_{vr}$ = virgin bulk density with the inert part removed [kg/m$^3$].
  \item $\rho_{cr}$ = char bulk density with the inert part removed [kg/m$^3$].
  \item $\mathrm{clip}$ = as in Eq.~(15).
\end{itemize}

\subsection*{Equation 21: linear char fraction}
\begin{equation}
\beta = \mathrm{clip}\!\left[
\frac{\rho_{v,\mathrm{bulk}} - \rho}{\rho_{v,\mathrm{bulk}} - \rho_{c,\mathrm{bulk}}},\, 0,\, 1\right]
\end{equation}
\textbf{What it is.} A second, simpler measure of charring, running the opposite
way from $\tau$. It equals $0$ in virgin material and $1$ in char. It drives
porosity and permeability, not the thermal properties. This is the degree-of-char
variable standard in the ablation literature.
\begin{itemize}
  \item $\beta$ = linear char fraction, dimensionless.
  \item $\rho_{c,\mathrm{bulk}}$ = bulk density of fully charred material
        [kg/m$^3$].
  \item $\rho$, $\rho_{v,\mathrm{bulk}}$ = as in Eqs.~(18), (14).
\end{itemize}
\textit{Note.} $\tau$ and $\beta$ are not interchangeable. They use different
formulas, run in opposite directions, and drive different physics.

\subsection*{Equation 22: porosity}
\begin{equation}
\varphi = \mathrm{clip}\left[\varphi_v + (\varphi_c - \varphi_v)\beta,\;
\varphi_{\min},\, \varphi_{\max}\right]
\end{equation}
\textbf{What it is.} How much void space is in the material, interpolated between
its virgin and char values as charring proceeds. Charred material is more porous.
\begin{itemize}
  \item $\varphi$ = porosity, dimensionless. First used in Eq.~(4).
  \item $\varphi_c$ = porosity of fully charred material.
  \item $\varphi_{\min}$, $\varphi_{\max}$ = the smaller and larger of $\varphi_v$
        and $\varphi_c$, used as clamps.
  \item $\varphi_v$, $\beta$ = as in Eqs.~(18), (21).
\end{itemize}

\subsection*{Equation 23: property mixing}
\begin{equation}
k_{\mathrm{mix}}(T,\tau) = \tau\, k_v(T) + (1-\tau)\, k_c(T),
\qquad
c_{p,\mathrm{mix}}(T,\tau) = \tau\, c_{p,v}(T) + (1-\tau)\, c_{p,c}(T)
\end{equation}
\textbf{What it is.} Blends the tabulated virgin and char property curves
according to char state. This is what makes the problem nonlinear, since the
properties depend on the solution.
\begin{itemize}
  \item $c_{p,v}(T)$ = tabulated specific heat of virgin material [J/kg/K].
  \item $c_{p,c}(T)$ = tabulated specific heat of char [J/kg/K].
  \item $k_{\mathrm{mix}}$, $c_{p,\mathrm{mix}}$, $k_v$, $k_c$, $\tau$, $T$ = as
        in Eqs.~(9), (12), (11), (20), (4).
\end{itemize}

\subsection*{Equation 24: formation enthalpy}
\begin{equation}
h_{s,\mathrm{form}}(\tau) = \tau\, e_v(T_{\mathrm{ref}}) + (1-\tau)\, e_c(T_{\mathrm{ref}})
\end{equation}
\textbf{What it is.} The baseline enthalpy of the material at the reference
temperature, blended by char state. The gap between the virgin and char values is
the heat of pyrolysis, which is why no separate parameter for it exists.
\begin{itemize}
  \item $h_{s,\mathrm{form}}$ = formation enthalpy [J/kg]. First used in Eq.~(8).
  \item $e_v(T)$ = integrated heat capacity of virgin material,
        $\int^T c_{p,v}\,dT'$ [J/kg].
  \item $e_c(T)$ = integrated heat capacity of char [J/kg].
  \item $T_{\mathrm{ref}}$ = reference temperature, $298.15$~K here.
  \item $\tau$ = as in Eq.~(20).
\end{itemize}

\subsection*{Equation 25: sensible enthalpy}
\begin{equation}
h_{s}(T,\tau) = \tau\, e_v(T) + (1-\tau)\, e_c(T) - h_{s,\mathrm{form}}(\tau)
\end{equation}
\textbf{What it is.} The energy stored in the solid by virtue of being hot,
measured from the reference temperature. Subtracting Eq.~(24) forces
$h_s = 0$ at $T_{\mathrm{ref}}$ for every char state, which is what isolates the
pyrolysis energy into Eq.~(8).
\begin{itemize}
  \item $h_s$ = sensible enthalpy [J/kg]. First used in Eq.~(4).
  \item All other symbols as in Eq.~(24).
\end{itemize}

\subsection*{Equation 26: emissivity}
\begin{equation}
\varepsilon = \tau\, \varepsilon_v + (1-\tau)\,\varepsilon_c
\end{equation}
\textbf{What it is.} How effectively the surface radiates heat away, blended by
char state in the same manner as the other properties.
\begin{itemize}
  \item $\varepsilon$ = emissivity, dimensionless, between $0$ and $1$.
  \item $\varepsilon_v$, $\varepsilon_c$ = emissivity of virgin material and char.
  \item $\tau$ = as in Eq.~(20).
\end{itemize}

\section{Pyrolysis-Gas Transport}

\subsection*{Equation 27: Darcy's law}
\begin{equation}
\dot m'' = -\frac{K_0 \rho_g}{\mu} \frac{\partial p_g}{\partial x}
\end{equation}
\textbf{What it is.} Gas flows through the porous char from high pressure to low
pressure, at a rate set by how permeable the material is and how viscous the gas
is. This replaces a momentum equation.
\begin{itemize}
  \item $K_0$ = permeability [m$^2$]. How easily gas passes through.
  \item $\mu$ = dynamic viscosity of the gas [Pa$\cdot$s].
  \item $p_g$ = gas pressure [Pa]. First used in Eq.~(4).
  \item $\dot m''$, $\rho_g$, $x$ = as in Eqs.~(9), (4).
\end{itemize}

\subsection*{Equation 28: discrete Darcy}
\begin{equation}
\dot m_i = -C_i \left(p_{g,i} - p_{g,i-1}\right),
\qquad
C_i = \frac{K_{0,i}\, \rho_g(T_{f,i}, p_{f,i})\, dA_i}
{\mu(T_{f,i}) \cdot \tfrac12\left(\Delta x_{i-1} + \Delta x_i\right)}
\end{equation}
\textbf{What it is.} Eq.~(27) written on the mesh. The bracketed group is bundled
into a single conductance $C_i$, so mass flux becomes conductance times pressure
difference. This is the same algebraic shape as Ohm's law.
\begin{itemize}
  \item $C_i$ = gas conductance of face $i$ [kg/s/Pa].
  \item $p_{f,i}$ = pressure interpolated to face $i$ [Pa].
  \item $\dot m_i$, $K_0$, $\mu$, $\rho_g$, $T_{f,i}$, $dA_i$, $\Delta x$ = as in
        Eqs.~(7), (27), (10), (1).
\end{itemize}

\subsection*{Equation 29: permeability blending}
\begin{equation}
K_0 = (1 - \beta_f) K_{0,v} + \beta_f K_{0,c}
\end{equation}
\textbf{What it is.} Char is more permeable than virgin material, so permeability
is interpolated as charring proceeds. Note this uses $\beta$, not $\tau$.
\begin{itemize}
  \item $K_{0,v}$, $K_{0,c}$ = permeability of virgin material and char [m$^2$].
  \item $\beta_f$ = linear char fraction evaluated at the face.
  \item $K_0$, $\beta$ = as in Eqs.~(27), (21).
\end{itemize}

\subsection*{Equation 30: gas mass conservation}
\begin{equation}
\frac{\partial (\varphi \rho_g)}{\partial t} + \frac{\partial \dot m''}{\partial x}
= \dot\omega_g
\end{equation}
\textbf{What it is.} Gas is conserved. The gas stored in the pores changes by
what flows in or out plus what pyrolysis creates.
\begin{itemize}
  \item All symbols previously defined. See Eqs.~(22), (4), (9), (18).
\end{itemize}

\subsection*{Equation 31: discrete gas conservation}
\begin{equation}
\dot m_{i+1} - \dot m_i = \dot\omega_{\mathrm{net},i}\, dV_i,
\qquad
\dot\omega_{\mathrm{net},i} = \dot\omega_{g,i}
- \frac{\left(\varphi \rho_g\right)^{n}_i - \left(\varphi \rho_g\right)^{n-1}_i}{\Delta t}
\end{equation}
\textbf{What it is.} Eq.~(30) on the mesh, rearranged so the storage term is
folded into the source. What remains is a steady balance, which is why the
pressure field is re-solved from scratch each step rather than marched in time.
\begin{itemize}
  \item $\dot\omega_{\mathrm{net},i}$ = net volumetric gas source after
        subtracting storage [kg/m$^3$/s].
  \item $n-1$ = the step before the current one.
  \item All other symbols previously defined.
\end{itemize}

\subsection*{Equation 32: direct march for the sealed-back case}
\begin{equation}
\dot m_1 = 0, \qquad
\dot m_{i+1} = \dot m_i + \dot\omega_{\mathrm{net},i}\, dV_i,
\qquad
p_{g,N+1} = p_{\mathrm{amb}}(t), \qquad
p_{g,i-1} = p_{g,i} + \frac{\dot m_i}{C_i}
\end{equation}
\textbf{What it is.} For the configuration used here, where the back face is
sealed and the front is open, Eq.~(31) can be solved by simple accumulation with
no matrix solve at all. Mass flux accumulates from the sealed back face, and
pressure is walked backward from the known ambient value at the surface.
\begin{itemize}
  \item $p_{\mathrm{amb}}(t)$ = ambient pressure at the surface [Pa],
        time-varying, taken from the aerothermal table.
  \item All other symbols previously defined.
\end{itemize}

\subsection*{Equation 33: gas properties}
\begin{equation}
\rho_g(T, p) = \frac{p}{R(T)\,\max(T, 250\,\mathrm{K})},
\qquad
c_{v,g} = c_{p,g}(T) - R(T)
\end{equation}
\textbf{What it is.} The ideal gas law, plus the standard relation between the
two gas specific heats. The gas composition changes with temperature, so the gas
constant is tabulated rather than fixed.
\begin{itemize}
  \item $R(T)$ = specific gas constant [J/kg/K], tabulated because the gas
        composition varies.
  \item $c_{p,g}(T)$ = specific heat of the gas at constant pressure [J/kg/K].
  \item $250$~K = a numerical floor, a guard inherited from the reference
        implementation.
  \item $\rho_g$, $c_{v,g}$, $p$, $T$ = as in Eqs.~(4), (12).
\end{itemize}

\section{Boundary Conditions}

\subsection*{Equation 34: surface energy balance}
\begin{equation}
q_w = \underbrace{\left(\rho_e u_e C_H\right) \Phi \left(h_r - h_w\right)}_{\text{convective}}
    + \underbrace{q_{\mathrm{rad}}}_{\text{incident}}
    + \underbrace{\frac{\dot m_{N+1}}{dA_{N+1}}\left(h_g^{\mathrm{abs}}(T_w) - h_w\right)}_{\text{pyrolysis gas}}
    - \underbrace{\varepsilon \sigma \left(T_w^4 - T_{\mathrm{ext}}^4\right)}_{\text{re-radiation}}
\end{equation}
\textbf{What it is.} The heat load at the exposed surface. Hot boundary-layer gas
heats the wall, incoming radiation adds to it, gas leaving the surface carries
energy across it, and the hot wall radiates energy back out. This is the term that
makes the aw2 case so much harder than aw1.
\begin{itemize}
  \item $q_w$ = net heat flux into the wall [W/m$^2$].
  \item $\rho_e u_e C_H$ = heat-transfer coefficient [kg/m$^2$/s]. A single
        tabulated quantity, not three separate factors to be multiplied. The
        subscript $e$ denotes boundary-layer edge conditions.
  \item $\Phi$ = blowing correction, dimensionless. Defined in Eq.~(35).
  \item $h_r$ = recovery enthalpy of the boundary-layer gas [J/kg].
  \item $h_w$ = wall enthalpy [J/kg]. Defined in Eq.~(36).
  \item $q_{\mathrm{rad}}$ = incident radiative flux [W/m$^2$], tabulated.
  \item $h_g^{\mathrm{abs}}(T_w)$ = absolute enthalpy of the pyrolysis gas at the
        wall [J/kg]. Distinct from $h_g$ in Eq.~(7), which is referenced to
        $T_{\mathrm{ref}}$.
  \item $T_w$ = wall temperature [K].
  \item $T_{\mathrm{ext}}$ = external temperature the surface radiates to [K].
  \item $\sigma$ = Stefan-Boltzmann constant, $5.670 \times 10^{-8}$
        W/m$^2$/K$^4$.
  \item $\varepsilon$, $\dot m$, $dA$ = as in Eqs.~(26), (7), (1).
\end{itemize}

\subsection*{Equation 35: blowing correction}
\begin{equation}
\varphi_b = \frac{2\lambda\, \dot m''_w}{\left(\rho_e u_e C_H\right)_0},
\qquad
\Phi = \frac{\varphi_b}{e^{\varphi_b} - 1}
\end{equation}
\textbf{What it is.} Gas leaving the surface pushes the boundary layer away from
the wall and reduces the convective heating. $\Phi$ is the factor by which the
heating is reduced, running from $1$ (no blowing) downward as blowing increases.
\begin{itemize}
  \item $\varphi_b$ = dimensionless blowing parameter. Note the subscript
        distinguishes it from porosity $\varphi$ in Eq.~(22). They are unrelated.
  \item $\lambda$ = tabulated blowing coefficient, dimensionless.
  \item $\dot m''_w$ = gas mass flux leaving the wall per unit area
        [kg/m$^2$/s].
  \item $\left(\rho_e u_e C_H\right)_0$ = the unblown heat-transfer coefficient,
        meaning the value before the correction is applied.
  \item $\Phi$ = as in Eq.~(34).
\end{itemize}

\subsection*{Equation 36: B-prime table lookup}
\begin{equation}
B'_g = \frac{\dot m_{N+1} / dA_{N+1}}{\left(\rho_e u_e C_H\right)_0 \Phi},
\qquad
h_w = h_w\!\left(T_w,\, B'_g\right)
\end{equation}
\textbf{What it is.} Surface chemistry, handled by table lookup rather than by
solving a chemistry problem at run time. $B'_g$ is a dimensionless blowing rate.
The wall enthalpy is read from a precomputed equilibrium table indexed by wall
temperature and $B'_g$.
\begin{itemize}
  \item $B'_g$ = dimensionless pyrolysis-gas blowing rate.
  \item All other symbols as in Eqs.~(34) and~(35).
\end{itemize}
\textit{Note.} Only one ambient-pressure table is read, so there is no pressure
interpolation of the surface chemistry.

\subsection*{Equation 37: surface ghost cell}
\begin{equation}
T_{N+1} = T_N + \frac{q_w}{k_{\mathrm{mix}}}\, \Delta x_N
\end{equation}
\textbf{What it is.} Imposes the boundary condition. The ghost-cell temperature
is set to whatever value makes the reconstructed gradient at the surface carry
exactly the heat flux $q_w$ from Eq.~(34).
\begin{itemize}
  \item $T_{N+1}$ = temperature of the ghost cell outside the heated surface [K].
  \item All other symbols previously defined.
\end{itemize}
\textit{Note.} For the gas-off case a prescribed wall temperature is imposed
instead, as $T_{N+1} = 2T_w(t) - T_N$.

\subsection*{Equation 38: back-face ghost cell}
\begin{equation}
T_0 = T_1 - \frac{q_{w,0}}{k_{\mathrm{mix}}}\, \Delta x_1, \qquad q_{w,0} = 0
\end{equation}
\textbf{What it is.} The back face is insulated, so no heat crosses it. With the
flux set to zero, the ghost cell simply copies the first interior cell.
\begin{itemize}
  \item $T_0$ = temperature of the ghost cell behind the back face [K].
  \item $q_{w,0}$ = heat flux at the back face [W/m$^2$], zero in all cases here.
  \item All other symbols previously defined.
\end{itemize}

\end{document}
